\chapter{Sistema de memoria y conocimiento}

\section{Propósito}

La memoria convierte ejecuciones aisladas en un sistema acumulativo. Sin ella,
cada ejecución vuelve a investigar los mismos paradigmas y puede repetir errores
ya detectados. Con ella, los resultados aceptados se transforman en conocimiento
recuperable.

La arquitectura combina cuatro tipos de almacenamiento:

\begin{figure}[H]
\centering
\begin{verbatim}
                 Artefacto aceptado
                         |
        +----------------+----------------+
        |                |                |
        v                v                v
      MinIO          PostgreSQL         Memoria
   bytes reales      ciclo de vida      KG + vectores
                                      Neo4j + Qdrant
\end{verbatim}
\caption{Separación de responsabilidades de almacenamiento.}
\end{figure}

\section{MemoryAgent}

MemoryAgent no es un agente conversacional libre. Es una tubería determinista
que usa modelos de lenguaje para extracción estructurada y juicio, pero mantiene
el control de flujo en código.

\begin{figure}[H]
\centering
\begin{verbatim}
MemoryAgent.run(stage, output, run_id)
  |
  +-- 1. Extraer nodos, relaciones y hechos
  +-- 2. Canonicalizar slugs de paradigmas
  +-- 3. Escribir grafo en Neo4j
  +-- 4. Indexar hechos en Qdrant
  +-- 5. Resolver duplicados y conflictos
  +-- 6. Persistir ciclo de vida en PostgreSQL
\end{verbatim}
\caption{Flujo interno de MemoryAgent.}
\end{figure}

La entrada de memoria siempre procede de artefactos aceptados. Esto protege el
sistema frente a borradores incorrectos y facilita explicar la memoria como
base de conocimiento curada.

\section{Modelo de extracción}

Cada etapa aceptada se convierte en una estructura validada:

\begin{lstlisting}[language=Python]
ExtractionResult(
    nodes=[NodeSpec(...), ...],
    relations=[RelationSpec(...), ...],
    facts=["atomic memory fact", ...],
    stage="researcher|formalizer|reasoner|builder",
    run_id="..."
)
\end{lstlisting}

La extracción usa salidas estructuradas y validación Pydantic. Las etiquetas y
slugs relevantes se comparan con vocabulario canónico. Si aparece un paradigma
nuevo, se usa un proceso de verificación antes de crear un nodo.

\section{Grafo de conocimiento}

Neo4j almacena identidad y relaciones científicas. Entre las etiquetas
principales se incluyen:

\begin{verbatim}
Paradigm, Paper, Author, Variable, Postulate,
Formulation, Parameter, Model, Reflection
\end{verbatim}

Y entre las relaciones:

\begin{verbatim}
SUPPORTS, CONTRADICTS, EXTENDS, MEASURES,
DERIVES_FROM, IMPLEMENTS, USES_VARIABLE,
HAS_PARAMETER, BELONGS_TO, CITES
\end{verbatim}

El grafo no almacena toda la verdad temporal. Las relaciones pueden enlazar a
una memoria mediante \texttt{memory\_id}, pero la confianza, validez temporal y
supersesión se mantienen en PostgreSQL.

\section{Memoria relacional}

La tabla \texttt{pipeline\_memories} conserva el ciclo de vida de los hechos.
Cada memoria incluye contenido, namespace, etapa de origen, importancia,
confianza, recuentos de corroboración y contradicción, accesos, intervalos de
validez y puntero de supersesión.

\begin{verbatim}
nuevo hecho
  -> no duplicado: crear memoria
  -> duplicado: omitir
  -> corroboracion: subir confianza
  -> contradiccion: bajar confianza y crear meta-memoria
  -> enriquecimiento: superseder memoria anterior
\end{verbatim}

Este diseño separa topología y ciclo de vida. Neo4j responde a preguntas de
relación; PostgreSQL decide qué conocimiento sigue vigente.

\section{Qdrant: recuperación densa y dispersa}

Qdrant contiene dos colecciones:

\begin{itemize}
\item \texttt{memories\_dense}: embeddings densos para similitud semántica.
\item \texttt{memories\_sparse}: BM25 nativo para coincidencia léxica.
\end{itemize}

El sistema no indexa fragmentos brutos de artefactos. Indexa hechos extraídos.
Los artefactos completos permanecen en MinIO. Esta decisión reduce ruido y
mantiene la recuperación enfocada en conocimiento aceptado.

\section{Herramienta de recuperación}

Los agentes acceden a la memoria mediante \texttt{retrieve\_knowledge}. La
herramienta combina varias señales:

\begin{figure}[H]
\centering
\begin{verbatim}
Consulta del agente
  -> reescritura de consulta
  -> recuperacion densa en Qdrant
  -> recuperacion BM25 en Qdrant
  -> recuperacion local en Neo4j
  -> fusion RRF
  -> reranking
  -> CRAG si la confianza es baja
  -> contexto formateado para el agente
\end{verbatim}
\caption{Pipeline de recuperación de conocimiento.}
\end{figure}

La recuperación de grafo empieza extrayendo entidades de la consulta. Después se
enlazan por coincidencia exacta o por índice vectorial de Neo4j y se recorren
vecinos locales. Se usa recorrido acotado para controlar coste y latencia.

\section{Corrección de recuperación}

Cuando los resultados recuperados son claramente fuertes, el sistema evita
trabajo adicional. Cuando son dudosos, CRAG clasifica su calidad:

\begin{verbatim}
todos correctos          -> usar contexto interno
mezcla correcta/erronea  -> conservar solo correctos
ambiguedad               -> complementar con web
todos incorrectos        -> fallback web
evaluador no disponible  -> devolver con advertencia
\end{verbatim}

Esta etapa evita que la memoria se convierta en una fuente de sesgo silencioso.
Es especialmente relevante en un sistema que almacena formulaciones y decisiones
técnicas de ejecuciones anteriores.

\section{Decisiones clave}

Las decisiones principales de la arquitectura de memoria son:

\begin{itemize}
\item Guardar solo artefactos aceptados tras revisión humana.
\item Mantener artefactos completos en MinIO y hechos recuperables en Qdrant.
\item Usar Neo4j para identidad relacional, no para ciclo de vida temporal.
\item Usar PostgreSQL como fuente de verdad para confianza y validez.
\item Combinar recuperación densa, BM25 y grafo en lugar de elegir un único
método.
\item Activar corrección solo cuando la confianza recuperada no es suficiente.
\end{itemize}

El resultado es una memoria más compleja que un vector store plano, pero más
adecuada para un sistema científico que debe justificar de dónde procede cada
afirmación.
